% =============================================================================
% manuscript.tex
% Elsevier `elsarticle` class. Author-year (Harvard) citation style via
% natbib, bundled with elsarticle-harv.bst.
%
% STATUS: skeleton manuscript. Sections 1 (Introduction), 2 (Background),
% 3 (Model structure and parameters), and 7 (Limitations and assumptions)
% are fully drafted from the repository's own documentation and code.
% Sections 4-6 and 8 are intentionally short stubs, marked with visible
% \todo{} markers, because this task was scoped to the first group only.
% See paper/README.md for the complete build instructions and the list of
% claims that still need a primary source.
%
% Table 1 (parameter table) was transcribed by reading vitd_model.py
% (module docstring and the ModelParams dataclass) on 2026-07-31. If the
% code changes, regenerate the table from the code; do not hand-edit values
% here without updating vitd_model.py first, per CLAUDE.md's single-source
% rule for parameter values.
% =============================================================================

\documentclass[preprint,12pt]{elsarticle}

\usepackage[T1]{fontenc}
\usepackage{amsmath,amssymb}
\usepackage{booktabs}
\usepackage{longtable}
\usepackage{array}
\usepackage{xcolor}
\usepackage[colorlinks=true,linkcolor=blue,citecolor=blue,urlcolor=blue]{hyperref}

% Visible, impossible-to-miss marker for anything the repository itself
% flags as unconfirmed, provisional, or not yet substantiated. Distinct from
% the plain LaTeX "%" comments used elsewhere in this file for source-only
% notes (build dates, transcription provenance) that a reader of the
% compiled PDF does not need to see.
\providecommand{\todo}[1]{\textcolor{red}{\textbf{[TODO verify: #1]}}}

\journal{}

\begin{document}

\begin{frontmatter}

\title{Scheduling, Not Just Dose: A Dynamic Photoadaptation Model for Cutaneous Vitamin D Synthesis Across Fitzpatrick Skin Types}

\author{Ahaan}
\address{\todo{affiliation not yet assigned -- institution, department, city, and country to be filled in before submission}}

\begin{abstract}
Vitamin D deficiency concentrates in people with darker skin at temperate latitudes, and existing models treat pigmentation as a fixed multiplier on a single-session dose-response curve. We present a three-state ordinary differential equation model of serum 25-hydroxyvitamin D, melanin index, and stratum corneum thickness, in which constitutive and facultative pigmentation, barrier change, and baseline-dependent feedback co-evolve with a specified ultraviolet schedule. Scheduling enters as an explicit decision variable. Every parameter carries an evidence-status tag: VERIFIED, PRIOR, SINGLE-SOURCE PRIOR, or UNCONFIRMED.
Four findings constrain what the model can claim. The homeostatic ceiling governing its scheduling behaviour has no confirmed literature origin. A widely used per-dose efficiency metric ranks schedules in near-exact opposition to health outcomes, with Spearman correlation near $-0.999$ against both mean 25(OH)D and time-integrated deficiency; holding cumulative dose fixed so that no schedule wins by dosing less, the optimal interval collapses to the finest spacing tested, for all fourteen limiting parameters across their documented ranges. We therefore report no interval recommendation. The melanin attenuation coefficient rests on one within-subject study of seven participants, whose cited 51\% attenuation reflects baseline feedback rather than pigment, and varying it moves steady-state output by under 1 nmol/L at realistic doses. The model also fails its only held-out benchmark, under-predicting the observed skin type V to skin type II dose ratio by roughly half.
Skin cancer risk appears as a qualitative ranking device and is deferred to future work. We offer a transparently flagged scheduling framework, not a clinical predictor.
\end{abstract}

\begin{keyword}
vitamin D \sep cutaneous photosynthesis \sep ultraviolet radiation \sep Fitzpatrick skin type \sep melanin \sep photoadaptation \sep stratum corneum \sep dose scheduling \sep homeostatic feedback
\end{keyword}

\end{frontmatter}

% =============================================================================
\section{Introduction}
% =============================================================================

Vitamin D deficiency concentrates in people with darker skin living at temperate latitudes. In the UK Biobank, profound deficiency (25-hydroxyvitamin D below 10 nmol/L) affects 0.3\% of White participants and 9.1\% of Asian participants, an adjusted odds ratio of 41.8 \citep{shraim2025}. The lowest quartile of ambient UVB exposure carries a 17-fold higher odds of profound deficiency than the highest quartile \citep{shraim2025}, and a companion analysis of the same cohort finds that ambient UVB, supplement use, and body-mass index all interact with ethnicity rather than acting additively, so a single population-average correction factor understates the gap for any one ethnic group \citep{brennan2024}. In a Manchester cohort of South Asian adults, almost all of whom were Fitzpatrick skin type V, fewer than 7\% reached sufficiency at any time of year, and winter medians sat near 5.8 ng/mL \citep{kift2013}. Farrar et al.\ quantified the gap under controlled conditions: at matched simulated-sunlight doses, skin type V participants needed roughly 2.5 times the UV exposure of white-skinned participants to achieve the same rise in serum 25(OH)D \citep{farrar2013}. A six-week head-to-head trial found the same pattern: most phototype II volunteers reached the 50 nmol/L sufficiency threshold, while over half the phototype V group reached only 25 nmol/L and none reached 50 \citep{felton2016}.

Scheduling changes outcomes independently of total dose. Low, sub-erythemal doses raise 25(OH)D more efficiently per unit dose than doses closer to the personal sunburn threshold, and splitting a UV budget across short, frequent sessions outperforms concentrating it into one long session, Webb et al.'s systematic review found \citep{webb2021}. Repeated narrowband UVB exposure plateaus in serum 25(OH)D after four sessions, with the next four adding almost nothing \citep{thieden2008}. Dose framing itself changes what the data appear to say about skin type: when dose is expressed as a fraction of each participant's own minimal erythemal dose (MED) rather than as a fixed physical dose, all six Fitzpatrick types produce statistically indistinguishable 25(OH)D gains per 20\% MED increment \citep{shih2018}. A model that reports a single number, expected rise per total standard erythemal dose (SED) delivered, captures none of this.

Existing dose-response models are cross-sectional or single-exposure fits. Webb et al.'s consolidated equation, $\Delta_{25(OH)D} = (9\pm0.5)\ln(\mathrm{SDD}) - (7.1\pm0.5)$, and Grigalavicius et al.'s hyperbolic alternative, $\Delta_{25(OH)D} = 54.8\,\mathrm{UV}/(22.9+\mathrm{UV})$, both predict one rise for one cumulative dose \citep{webb2021,grigalavicius2015}. Neither updates a participant's pigmentation state as sessions accumulate; both hold Fitzpatrick type, or a fixed melanin index, constant for the whole exposure period. That simplification is reasonable for a single clinical exposure. It breaks down over a multi-week course, where each session changes the melanin and stratum corneum state that governs the next session's synthesis. A darker-skinned participant on a three-times-weekly schedule is not the same optical system in week 1 and week 9, and the published literature compares endpoints across such courses without ever modeling the trajectory between them.

This paper describes a model built to close that gap: three coupled state variables, serum 25(OH)D, melanin index, and stratum corneum thickness, driven by a caller-specified UV pulse schedule, so that the interval between sessions becomes a first-class model output alongside their total. Section~\ref{sec:model} lays out the governing equations and the full parameter set. Rather than optimizing one scalar dose-response slope, the accompanying codebase poses the scheduling question the way a person or clinician actually asks it: given a fixed number of sessions over a fixed horizon, when should they fall; given a target of staying above a deficiency threshold, what schedule costs the least in skin-cancer risk; and given an ambient UV index and skin type, how many minutes of sun are enough.

Several of this model's central results are provisional, and we state them here rather than only in Section~\ref{sec:limitations}. The homeostatic ceiling that governs feedback inhibition has no confirmed literature source, and ablation shows it, not pigmentation, drives the model's ranking of exposure intervals. A per-dose efficiency metric used earlier in this project's development rewards schedules that end below the participant's starting vitamin D level and ranks schedules in near-exact opposition to the health outcomes that actually matter; a net-change alternative was reported to fix this, but it shares the same ratio structure, and under a framing immune to that structure, cumulative UV dose held fixed so no schedule can win by dosing less, the model's optimal interval collapses to the finest spacing tested across every parameter combination checked. No interval recommendation should currently be read out of this model. The melanin attenuation coefficient traces to one seven-person study whose most-quoted number turns out to be mostly a baseline-feedback effect, and melanin attenuation barely moves the model's own output at realistic doses. The model under-predicts the real-world skin-type dose penalty against the one benchmark it was never calibrated to. Skin-cancer risk is represented only as an order-of-magnitude ranking device, deliberately not developed further in this paper. Section~\ref{sec:limitations} carries the complete register of these and other open issues, cross-referenced to the repository's own tracking documents.

% =============================================================================
\section{Background}
\label{sec:background}
% =============================================================================

\subsection{Cutaneous vitamin D photosynthesis and its dose-response form}

UVB radiation, roughly 280 to 315 nm, converts 7-dehydrocholesterol in the epidermis to previtamin D3, which isomerizes thermally to D3 and is then hydroxylated in the liver to 25-hydroxyvitamin D, the standard clinical biomarker of vitamin D status. Webb et al.'s systematic review of 22 artificial-source and 14 solar-exposure studies fit a logarithmic dose-response form to the pooled data: $\Delta_{25(OH)D} = 9.51\ln(\mathrm{SDD}) - 7.6$ on seasonally adjusted data ($r^2 = 0.66$), consolidated across their equations to $(9\pm0.5)\ln(\mathrm{SDD}) - (7.1\pm0.5)$ \citep{webb2021}. Here SDD, the standard vitamin-D dose, is 100 J/m\textsuperscript{2} of vitamin-D-effective radiation, an analogue of the SED used for erythema. Grigalavicius et al.'s independent meta-regression of 18 studies (699 subjects) instead fit a saturating hyperbola, $\Delta_{25(OH)D} = 54.8\,\mathrm{UV}/(22.9+\mathrm{UV})$, $r = 0.78$, with a plateau near 55 nmol/L and half-maximum at 23 SED \citep{grigalavicius2015}. The two forms are not identical, but they agree over the sub-erythemal dose range that matters clinically, and both describe diminishing per-SED returns as cumulative dose rises. The model presented here adopts Webb's logarithmic form directly in its per-session synthesis term.

Source spectrum changes the conversion between physical dose and vitamin-D-effective dose. Webb et al.\ tabulate vitamin-D-effective to erythema-effective (VDE:EE) ratios of 2.24 for narrowband TL-01, 1.25 for broadband TL-12, 1.55 for UV6, and 1.89 for summer sun \citep{webb2021}; Grigalavicius et al.\ report comparable but not identical values from an independent set of studies \citep{grigalavicius2015}. Young et al.'s in vivo test of the CIE previtamin-D3 action spectrum found it did not fit measured serum 25(OH)D3 responses across five different UV sources unless shifted roughly 5 nm toward shorter wavelengths \citep{young2021pnas}. Osmancevic et al.\ found that broadband UVB phototherapy raised 25(OH)D3 more than narrowband UVB in the same psoriasis population, despite narrowband being the more erythema-efficient source \citep{osmancevic2009}, and a synthesis of three phototherapy modalities from the same group placed heliotherapy, broadband UVB, and narrowband UVB in that same order \citep{osmancevic2010}: source-specific conversion factors cannot collapse to one constant. The model implements this as a per-study calibration table (\texttt{SOURCE\_CALIBRATION\_TABLE}) rather than a single multiplier, with a default corresponding to broadband solar exposure.

A single session cannot substitute for several small ones, because previtamin D3 production saturates within a session. \citet{holick1981} \todo{cited in this repository only as ``Holick 1981, Science 211:590--593''; full author list not recorded} reported that continued UV exposure drives previtamin D3 to a plateau at roughly 10 to 15\% of the original 7-dehydrocholesterol content, after which further UV photoisomerizes it to biologically inert lumisterol3 and tachysterol3, and that pigmentation and latitude change the time to reach this plateau without changing its height. This model implements that mechanism as a session-averaged saturating factor on synthesis, ablatable and not enabled by default (Section~\ref{sec:model}), and it is the reason a genuine benefit from splitting a fixed UV budget across sessions can exist independently of the feedback-driven artifact discussed in Section~\ref{sec:limitations}.

\subsection{Melanin and constitutive pigmentation as an optical filter}

Melanin absorbs UVB before it reaches the 7-dehydrocholesterol pool, and the size of that effect is contested. Between Fitzpatrick skin types II and VI, the measured penalty is only 1.3- to 1.4-fold under matched-dose lamp exposure (Arimed B: 1.29; TL-01: 1.43), small enough that Young et al.\ titled the resulting paper accordingly: melanin has a small inhibitory effect \citep{young2020}. Bogh et al.\ went further, finding no significant relationship between either constitutive or facultative pigmentation and the UVB-induced rise in 25(OH)D, even when comparing matched pairs of dark-skinned (types V to VI) and fair-skinned (types I to IV) participants whose pigmentation differed sharply ($P=0.008$) but whose response did not ($P=0.7$) \citep{bogh2010jid}. The same pattern holds at the individual level: measured skin pigmentation explained only 1.9 to 2.6\% of individual variance in UVB-induced 25(OH)D response across two separate cohorts, pigment-associated single-nucleotide polymorphisms explained far more, and identical UVB exposure produced final serum levels ranging from 73 to 216 nmol/L across participants \citep{datta2019,datta2016}.

Naturalistic and whole-body studies report a larger gap. Farrar et al.\ measured a roughly 2.5-fold dose penalty for skin type V \citep{farrar2013}, and a matched six-week trial found phototype V participants gained half the 25(OH)D of phototype II participants under an identical protocol \citep{felton2016}. Kift et al.\ add a behavioral layer on top of any pigmentation effect: South Asian participants carried barely half the personal UV dosimeter readings of white participants despite similar time outdoors \citep{kift2013}. Bogh's null result used a UVB source (TL-12) with substantial output below 300 nm, wavelengths that penetrate less deeply and may largely bypass the melanin layer; Young's result used solar-simulated sources closer to natural sunlight. That spectral difference does not resolve the gap between a roughly 1.3-fold lamp-measured penalty and a 2.5- to 3-fold naturalistic one, and Section~\ref{sec:limitations} returns to it: this model's constitutive-pigmentation coefficient is calibrated to the smaller, lamp-derived ratio, so every real-world minutes-of-sun estimate it produces for darker skin is likely optimistic.

\subsection{Facultative photoadaptation: tanning and stratum corneum change}

Constitutive pigmentation is the melanin present without recent sun exposure; facultative pigmentation is the additional melanin that UV exposure induces, and it fades once exposure stops. \citet{markiewicz2020} caution that visible tanning often reflects redistribution and oxidation of existing melanin rather than large new synthesis, a mechanistic detail this model does not represent and that would, if material, argue against treating tanning as a simple saturating induction process. Ravnbak et al.\ measured tanning kinetics across skin types II to V and two ethnic groups, finding that repeated exposure lowers the dose needed to reach a given pigmentation endpoint by a factor of roughly 2 after six exposures and 3 after twelve, and that full fading takes four to five months regardless of the exposure frequency that produced the tan \citep{ravnbak2008}; a companion UVA study found five to six months \citep{ravnbak2009}. This model uses the Ravnbak fading data for its melanin decay rate; an earlier version of the project had cited a different paper (Tadokoro et al., 2005) for this parameter, which measures pigmentation only over the first week post-exposure and reports no decay timescale at all, and that citation has been corrected.

De Boer et al.\ ran a nine-week repeated low-dose UVR protocol (0.8 SED, three times weekly) and reported the most complete photoadaptation dataset available to this project: MED rose 58\% in phototypes I to III and 33\% in phototypes IV to VI; melanin index rose 15\% and 9\% respectively; stratum corneum thickness changed 23\% and 36\% respectively, but not monotonically, falling from 5.8 to 4.3 \textmu m (about 26\% below baseline) at five weeks before rising to 7.5 \textmu m by nine weeks \citep{deboer2025}. In the same dataset, epidermal thickness barely moved, 108.6 to 109.4 \textmu m ($P=0.718$), and the authors attribute the stratum corneum change to edema rather than acanthosis (true corneocyte-layer thickening) \citep{deboer2025}. That distinction matters for any model, including this one, that treats stratum corneum thickness as added optical path length attenuating UV: fluid retention is not the same physical mechanism as more corneocyte layers, and Section~\ref{sec:limitations} flags this as a premise that may be wrong even where the fitted numbers happen to work. Pearse et al.\ confirm the broader phenomenon mechanistically, an 87\% increase in stratum corneum thickness after six weeks of UVB at twice the minimal erythemal dose, but report only cumulative six-week endpoints, not the rate constants a dynamic model needs \citep{pearse1987}. A summer-exposure cohort of 258 Chinese participants adds a further caveat this model does not represent: barrier changes after natural sun exposure were sex-dependent, with basal transepidermal water loss rising with dose in men but falling with dose in women \citep{liu2010}.

\subsection{Homeostatic feedback and catabolism}

Vitamin D response depends on where a participant starts. Bogh et al.\ fit $\Delta_{25(OH)D} = 7.1 - 0.223 \times \mathrm{baseline} + 11.0\ln(\mathrm{UVB\ dose}) + 5.5\ln(\mathrm{BSA\%})$ across 92 participants ($R^2=0.47$), establishing that higher baseline 25(OH)D predicts a smaller rise for the same dose \citep{bogh2011bjd}. Karppinen et al.\ found something sharper than diminishing returns: under identical vernal solar exposure to just 7\% of body surface area, participants who started below 90 nmol/L rose by 6.0 nmol/L on average, while participants who started above 90 nmol/L fell by 6.7 nmol/L \citep{karppinen2017}. That is active downregulation, not merely a saturating response curve. Datta et al.\ showed the same baseline dependence in the decay direction: post-UVB half-life ranged from 89 days at a mean starting level of 132 nmol/L to 199 days at 43 nmol/L \citep{datta2017}; this model fits a continuous power-law decay rate to those three group means (state-dependent clearance, $t_{1/2}(D)=2953\,D^{-0.72}$ days) rather than assuming a single fixed half-life. A seven-day narrowband UVB course in Finnish winter reached a peak two weeks after the last exposure and remained above baseline two months later \citep{vahavihu2010a}, consistent with a slow, state-dependent decay rather than a short plasma half-life; a separate direct comparison found six narrowband exposures (7.8 SED cumulative) raised 25(OH)D by 11.4 nmol/L against 4.1 nmol/L for three weeks of 400~IU daily oral D3 in the same season \citep{alahouhala2012}.

Physiologically, this regulation is understood to act on catabolism: CYP24A1 hydroxylates both 25(OH)D and its active metabolite 1,25(OH)\textsubscript{2}D toward inactivation, and loss-of-function mutations in this enzyme produce sustained, unregulated elevation of the active hormone \citep{carpenter2017}. The model represents this feedback as a linear inhibition term, $\max(0,\,1-D/D_{\max})$, multiplying every session's synthesis, a simplification of that catabolic mechanism rather than a direct implementation of it. $D_{\max}$ carries the largest influence of any parameter in the model and the weakest evidentiary support. Its origin is unconfirmed: an earlier project revision attributed it to Bogh et al.'s 2010 study, but that paper's reported ceiling was 116 nmol/L, and the PMID once cited for it belongs to an unrelated obesity paper (\citealt{doumatey2010}; see \citealt{bogh2010jid} for the actual Bogh 2010 study). \todo{$D_{\max}=150$ nmol/L origin is unconfirmed. No primary source in this repository supports this value; see Section~\ref{sec:limitations} for why this matters more than any pigmentation parameter.} A more physiologically grounded alternative treats CYP24A1-mediated catabolism and 25-hydroxylation efficiency, not a hard synthesis ceiling, as the regulated step: a 2025 Communications Medicine study reports 25-hydroxylation efficiency falling markedly above an inflection near 56 nmol/L \citep{commsmed2025}. \todo{Full bibliographic details (authors, article title, volume, pages) for the Communications Medicine 2025 source are not recorded anywhere in this repository; only the PMC identifier (PMC12894941) is available. The model implements this biphasic form as an ablatable alternative, not yet the default, and the steepness of its transition is an unconfirmed guess: the source reports only the inflection location, not the curve's shape.}

\subsection{Dose framing: absolute SED versus fractional MED}

Two dose framings produce apparently contradictory literatures, and the contradiction is an artifact of framing rather than a disagreement about biology. Shih et al.\ dosed skin types I through VI at 20\%, 40\%, 60\%, and 80\% of each participant's own MED and found equivalent 25(OH)D gains per 20\% MED increment across all six types, 1.6 nmol/L on average (95\% CI 1.1 to 2.0) \citep{shih2018}. Under this fractional-MED framing, the apparent skin-type effect on vitamin D response largely disappears. Farrar et al., Felton et al., and Kift et al., by contrast, all used a fixed absolute UV dose (or measured naturalistic exposure) across skin types, and all three found that darker-skinned participants remained under-dosed relative to sufficiency targets that lighter-skinned participants reached \citep{farrar2013,felton2016,kift2013}. A fixed absolute dose asks less, biologically, of a person with more melanin, since the same physical dose represents a smaller fraction of their much higher sunburn threshold; a fixed fractional-MED dose does not. Neither framing is wrong. A model, or a public-health recommendation, that does not state which one it is using will look self-contradictory the moment it is compared against a study using the other. This model requires every simulated scenario to declare its dose framing explicitly (fixed absolute SED, fixed cumulative SED, fixed per-session SED, or fixed fractional MED) for exactly this reason, and results generated under one framing should never be read across to a claim made under another.

\subsection{UV exposure and skin cancer risk as the competing objective}

Any UV exposure recommendation trades vitamin D against skin cancer risk, and the epidemiology splits by cancer type in a way that matters for scheduling. Squamous cell carcinoma tracks cumulative lifetime UV dose. Basal cell carcinoma and melanoma track intermittent, intense exposure and sunburn history more than cumulative dose: Elwood and Jopson's meta-analysis of 29 case-control studies reported an odds ratio of 1.71 for intermittent exposure and, in the same analysis, an odds ratio of 0.86, protective, for heavy occupational (continuous) exposure \citep{elwood1997}. \todo{Elwood \& Jopson 1997 is cited in this repository only by author and year, with the specific odds ratios attributed to it; journal, volume, and page numbers are not recorded and need verification against the primary source before this figure is quoted elsewhere.} Leiter and Calzavara-Pinton report a consistent pattern, basal cell carcinoma associated with intermittent exposure (especially in childhood) and squamous cell carcinoma with chronic, cumulative exposure \citep{leiter2009,calzavarapinton2015}, an ordering an early case-control analysis and a large modern occupational cohort both support \citep{armstrong1988,osullivan2026}. \todo{Leiter 2009, Calzavara-Pinton 2015, Armstrong 1988, and O'Sullivan 2026 are recorded in this repository only as author-year fragments (the O'Sullivan cohort is described as approximately 2.4 million Ontario workers, low-intermittent hazard ratio 1.26 versus high-continuous hazard ratio 1.13); full bibliographic detail is not available for any of the four and all need verification before submission.} The practical implication runs counter to naive dose minimization: for a fixed total UV budget, spreading it across many small sub-erythemal sessions can lower melanoma and basal cell carcinoma risk relative to concentrating the same dose into fewer, larger sessions, even though cumulative dose, and therefore squamous cell carcinoma risk, is unchanged.

We deliberately do not develop this into a quantitative risk model in the present work. The accompanying code implements only a qualitative, order-of-magnitude ranking device: cumulative dose scales a squamous-cell-type index linearly, per-session dose relative to a skin type's MED scales an intermittency-type index cubically, and a susceptibility factor falls with Fitzpatrick type to match the roughly tenfold gap in incidence between very fair and deeply pigmented populations. None of these three functional forms is fitted to incidence data; all three are structural choices consistent with the epidemiology's qualitative direction, not calibrated magnitudes. \todo{The papers that would let this become a quantitative model, Eadie 2021 on UV dose-response for skin cancer \citep{eadie2021}, Armstrong and Kricker 2001 on intermittent versus cumulative exposure \citep{armstrongkricker2001}, and a 2023 PLOS ONE study on intermittent versus continuous UV and melanoma \citep{plosone2023}, are not present in this repository and have not been obtained. Do not cite the cancer-risk output of this model as a validated absolute-risk estimate; it is a same-total-dose ranking device only, and only a same-total-dose ranking device.} We flag this explicitly as future work and return to it in Section~\ref{sec:limitations}.

% =============================================================================
\section{Model structure and parameters}
\label{sec:model}
% =============================================================================

\subsection{Governing equations}

The model tracks three coupled states over time $t$ in days: serum 25(OH)D, $D(t)$, in nmol/L; an effective melanin index, $M(t)$, on a dimensionless 0 to 1 scale anchored to the Datta et al.\ reflectance convention (0 corresponds to Fitzpatrick type I, 1 to untanned Fitzpatrick type VI); and additional stratum corneum thickness above baseline, $SC(t)$, dimensionless. Between UV pulses, the states evolve as

\begin{align}
\frac{dD}{dt} &= S\bigl(UV(t), M(t), SC(t), D(t), F, \mathrm{bsa}\bigr) - k_{\mathrm{decay},D}(D)\, D, \\
\frac{dM}{dt} &= \alpha\, UV(t)\,(M_{\max} - M) - \beta\,(M - M_{\mathrm{baseline}}(F)), \\
\frac{dSC}{dt} &= \alpha_{SC}\, UV(t)\,(SC_{\max} - SC) - k_{\mathrm{decay},SC}\, SC - k_{\mathrm{damage}}\, UV(t)\, SC,
\end{align}

\noindent where $F$ is Fitzpatrick skin type (1 to 6) and $\mathrm{bsa}$ is the fraction of body surface area exposed. Each UV pulse of dose $\mathrm{sed}$ (in SED) delivers a synthesis increment

\begin{align}
S(\mathrm{sed}, M, SC, D, F, \mathrm{bsa}) = {}& A\,\ln(1+\mathrm{sed}\cdot E_{\mathrm{source}})\cdot \mathrm{bsa\_factor}(\mathrm{bsa}) \nonumber\\
&\times \exp\bigl(-K_{\mathrm{const}}\, M_{\mathrm{baseline}}(F)\bigr) \times \exp\bigl(-K_{\mathrm{ATTEN}}\max(0,\,M-M_{\mathrm{baseline}}(F))\bigr) \nonumber\\
&\times \exp(-K_{SC}\, SC) \times \max\!\left(0,\, 1-\frac{D}{D_{\max}}\right),
\end{align}

\noindent with $\mathrm{bsa\_factor}(\mathrm{bsa}) = (\mathrm{bsa}/\mathrm{BSA_{REF}})^{\mathrm{BSA\_EXP}}$, a sub-linear scaling that equals 1 at the reference exposure fraction. The melanin induction rate $\alpha$ itself depends on session dose through a Hill-type saturation, $\alpha(\mathrm{sed}) = \alpha_{\mathrm{base}}(F) + (\alpha_{\max}-\alpha_{\mathrm{base}}(F))\,\mathrm{sed}^2/(\mathrm{sed}^2+UV_{\mathrm{halfmax}}(F)^2)$, so a single large session induces more tanning per unit dose than several small ones. The linear feedback term $\max(0,1-D/D_{\max})$ is the model's default; a biphasic alternative grounded in 25-hydroxylation efficiency (Section~\ref{sec:background}) is implemented but not enabled by default.

Between pulses, all three states decay or regress toward their unstimulated values, with $D$ following the state-dependent clearance $k_{\mathrm{decay},D}(D) = \ln(2)/\bigl(2953\,D^{-0.72}\bigr)$ fitted to \citet{datta2017}'s three group half-lives. The implementation used for simulation additionally splits $SC$ into two separately evolving sub-processes, a slow hyperplasia-type build-up and a fast UV-driven acute-damage term, with $SC$ equal to their difference; this lets the net stratum corneum state fall below its zero baseline, which \citet{deboer2025}'s data show it does at five weeks. The acute-damage rate constant, $k_{\mathrm{damage}}$, remains fixed at zero throughout this manuscript: it is not identifiable from the three group-median timepoints available (see Section~\ref{sec:limitations}).

\subsection{Objectives}

Rather than optimizing a single scalar dose-response slope, the accompanying codebase (\texttt{vitd\_objectives.py}) poses three decision problems directly. Given a fixed number of sessions over a fixed horizon, which spacing maximizes mean, final, or minimum 25(OH)D, or minimizes time-integrated deficiency below a target; because session count and per-session dose are both fixed, cumulative UV dose is identical across every candidate spacing, so no schedule can win the comparison simply by delivering less UV. Given a deficiency target and a horizon, which combination of dose-per-session and inter-session interval traces the best available trade-off between time spent below the target and the two skin-cancer risk indices from Section~\ref{sec:background}, reported as a Pareto frontier rather than collapsed into one weighted score. Given an ambient UV index, skin type, and exposed body-surface-area fraction, how many minutes per session reach a serum 25(OH)D target, and does that exposure cross the erythema threshold for that skin type.

This framing choice is not cosmetic, and a diagnostic sweep shows why. \texttt{scripts/parameter\_optima.py} computed the optimal inter-session interval under two framings side by side, for fourteen parameters that penalize synthesis or burden the body (constitutive and facultative pigmentation, stratum corneum change, homeostatic feedback, and exposure geometry), each swept independently across its full documented prior range, and jointly by 200-draw Monte Carlo. Under a fixed cumulative UV budget, split across sessions so that no schedule can win by dosing less, the optimal interval is one day, the finest spacing tested, for every one of the fourteen parameters at every tested value; the joint Monte Carlo gives a 90\% interval of exactly $[1,1]$ days. Under a per-SED ratio objective (net 25(OH)D change per unit dose delivered), where cumulative dose is free to vary across candidate schedules, a multi-day optimum appears instead, and only two of the fourteen parameters, exposed body-surface-area fraction and the synthesis coefficient $A$, both dose-magnitude parameters rather than photoadaptation parameters, move that per-SED optimum at all. Section~\ref{sec:limitations} discusses what this means for any interval recommendation quoted from this project, including from earlier drafts of this manuscript.

\subsection{Parameters}

Table~\ref{tab:params} lists every scientific parameter used in the model, transcribed from \texttt{vitd\_model.py} (module docstring and \texttt{ModelParams} dataclass), \texttt{vitd\_physiology.py}, and \texttt{vitd\_objectives.py} on 2026-07-31. The evidence-status column uses the same four tags the source repository uses internally: \textbf{VERIFIED} (the value or its functional form is directly supported by a specific measurement in the cited source), \textbf{PRIOR} (a plausible, literature-informed value that has not been independently fitted or replicated), \textbf{SINGLE-SOURCE PRIOR} (supported by exactly one study, with no independent replication located), and \textbf{UNCONFIRMED} (no primary source has been located for this value, or the source's own premise is in question). Several rows carry additional caveats beyond the tag; those caveats are load-bearing and should travel with the number whenever it is quoted elsewhere.

{\footnotesize
\begin{longtable}{@{}p{2.1cm}p{3.6cm}p{1.6cm}p{2.6cm}p{2.0cm}p{3.4cm}@{}}
\caption{Model parameters, transcribed from \texttt{vitd\_model.py} on 2026-07-31. Do not restate these values elsewhere without re-checking them against the source file; parameter values live in exactly one place in the source repository.}
\label{tab:params}\\
\toprule
Symbol & Meaning & Units & Value & Evidence status & Citation / notes \\
\midrule
\endfirsthead
\multicolumn{6}{c}{\tablename\ \thetable\ (continued)}\\
\toprule
Symbol & Meaning & Units & Value & Evidence status & Citation / notes \\
\midrule
\endhead
\bottomrule
\endfoot

$A$ & Synthesis scaling coefficient & nmol/L per $\ln$-unit & 2.353 & PRIOR & Calibrated to reproduce a 1.27 nmol/L rise per 1 SED in untanned type III skin at reference BSA; order-of-magnitude consistent with \citet{young2021pnas}'s 2.0--5.0 nmol/L/SED fair-skin slopes at 85\% BSA. \citet{borecka2024} and \citet{chalcraft2020} provide partial low-dose response data in phototype I--III consistent in direction but not independently used to set this value. \\

$K_{\mathrm{const}}$ & Constitutive (baseline) pigmentation inhibition & dimensionless & 0.33 & VERIFIED & $\ln(1.36)/(M_{VI}-M_{II})$ from \citet{young2020}'s type II:VI ratio; directionally supported by \citet{libon2013} and \citet{armas2007}. \\

$K_{\mathrm{ATTEN}}$ & Facultative melanin attenuation coefficient & dimensionless & 0.37 & SINGLE-SOURCE PRIOR & Derived from \citet{matsuoka1990jlcm} (abstract only, $n=7$, group means); \citet{young2020} independently implies a small effect of similar order. \todo{No independent replication of the Matsuoka within-subject tanning-attenuation result exists.} \\

$D_{\max}$ & Homeostatic feedback ceiling (linear mode) & nmol/L & 150.0 & UNCONFIRMED & \todo{Origin not located in any source in this repository. Drives the model's interval-ranking result more than any other parameter (Section~\ref{sec:limitations}).} \\

$\alpha$ (by FST) & Melanin induction base rate & 1/SED & 0.0025 (I--III), 0.0020 (IV--VI) & PRIOR & Original default calibrated against light-skin photoadaptation kinetics \citep{miyamura2011}. \todo{Miyamura 2011 is cited in this repository's gap register only as ``Miyamura 2011''; full bibliographic detail not recorded.} Direction supported by \citet{deboer2025} (FST-stratified MI gains). A dedicated re-fit against \citet{young2020}'s trial data found $\alpha$ structurally unidentifiable from that dataset (92.5\% of fits hit search bounds). \\

$\alpha_{\max}$ & Melanin induction rate at high per-session dose & 1/SED & 0.02 & PRIOR & Hill-form ceiling; not independently sourced. \\

$UV_{\mathrm{halfmax}}$ (by FST) & Hill half-max dose for melanin induction & SED & 2.0 (I--III), 5.0 (IV), 8.0 (V), 10.0 (VI) & PRIOR & Aligned to \citet{ravnbak2008} Table IV post-adaptation multi-exposure minimal melanogenic dose, not the single-exposure MMD in the same paper's Table I. \\

$\beta$ & Facultative melanin decay rate & 1/day & $\ln 2/56$ (56 d half-life) & VERIFIED & 56 d half-life fits the 4--6 month total fade reported by \citet{ravnbak2008,ravnbak2009}. Superseded citation: this value was previously attributed to Tadokoro et al.\ (2005), a 0--7 day study reporting no decay timescale; corrected. \\

$M_{\mathrm{baseline}}$ (by FST) & Constitutive melanin index & dimensionless, 0--1 & I: 0.00, II: 0.07, III: 0.21, IV: 0.44, V: 0.78, VI: 1.00 & PRIOR & Anchored to ITA (individual typology angle) unexposed-skin means from an Arimed~B trial dataset ($n=44$). Types V and VI have no direct published melanin-index validation and are extrapolated. \todo{\citet{datta2019} shows measured pigmentation explains only 1.9--2.6\% of individual response variance; treat this row as a group-level phenotype marker, not an individual predictor.} \\

$M_{\mathrm{TAN\_HEADROOM}}$ & Facultative tanning headroom above baseline & dimensionless & 0.20 & PRIOR & Magnitude not separately identified; \citet{ravnbak2008} reports facultative gains on an instrument-specific scale not directly convertible. \\

$k_{\mathrm{decay},D}(D)$ & State-dependent 25(OH)D clearance & 1/day; $t_{1/2}=2953\,D^{-0.72}$ d & function of $D$ & VERIFIED & Power law fitted to \citet{datta2017}'s three group half-lives (89 d at 132 nmol/L, 148 d at 65 nmol/L, 199 d at 43 nmol/L); the power-law form itself is not from the source paper. \\

$SC_{\max}$ (by FST) & Stratum corneum thickening ceiling & dimensionless & 0.23 (I--III), 0.36 (IV--VI) & VERIFIED & FST-stratified from \citet{deboer2025} Supplementary Table S1 (+23\% FP I--III, +36\% FP IV--VI over 9 weeks); anchored to only two informative timepoints per stratum. \\

$\alpha_{SC}$ & Stratum corneum hyperplasia rate & 1/SED & 0.12 & PRIOR & Fitted to reach approximately 90\% of $SC_{\max}$ within \citet{deboer2025}'s 9-week, 3x/week protocol; fitting uncertainty from three group-mean timepoints only. \\

$k_{\mathrm{decay},SC}$ & Stratum corneum regression rate & 1/day & $\ln 2/28$ (28 d half-life) & PRIOR & Generic epidermal-turnover literature value, not fitted specifically to the cited dataset. \\

$K_{SC}$ & Stratum corneum optical attenuation coefficient & dimensionless & 1.0 & UNCONFIRMED & \todo{In the same dataset used to calibrate this parameter, epidermal thickness barely changed (108.6 to 109.4 \textmu m, $P=0.718$) while stratum corneum thickness swung $\pm26\%$; \citet{deboer2025} attribute the change to edema, not acanthosis. The optical-attenuation interpretation of this parameter may rest on a mechanism that is not actually occurring, independent of whether the fitted number happens to improve model fit.} \\

$k_{\mathrm{damage}}$ & UV-driven acute stratum corneum barrier disruption & 1/SED & 0.0 (pinned) & UNCONFIRMED & \todo{Not identifiable from the three group-median timepoints in \citet{deboer2025}; the best achievable fit misses 71--75\% of signal amplitude. Held at zero rather than fitted to a value that overstates the available evidence.} \\

$k_{\mathrm{repair},SC}$ & Acute stratum corneum damage repair rate & 1/day & $\ln 2/7$ (7 d half-life) & PRIOR & Not separately measured in any cited source; chosen so a modeled dip recovers within one protocol cycle. \\

$E_{\mathrm{source}}$ (solar broadband, default) & Source/spectrum vitamin-D-effective dose weighting & dimensionless & 1.00 (reference) & PRIOR & See Table~\ref{tab:sources} for the full per-source calibration table. \\

$\mathrm{BSA_{REF}}$ & Reference exposed body-surface-area fraction & fraction & 0.35 & VERIFIED & From \citet{farrar2013}'s protocol (35\% BSA exposed); \citet{bogh2011bjd} tested only 6--24\% BSA. \\

$\mathrm{BSA\_EXP}$ & Sub-linear body-surface-area scaling exponent & dimensionless & 0.50 & PRIOR & Consistent with \citet{bogh2011bjd}'s $\ln$(BSA\%) term and the plateau above 33\% BSA reported by \citet{matsuoka1990jaad}. \\

$D_{0,\mathrm{default}}$ & Default simulation starting serum 25(OH)D & nmol/L & 25.0 & Model convention & Simulation initial condition, not a literature-derived parameter. \\

$D_{\mathrm{half,feedback}}$ & Saturating-feedback half-decay level (non-default mode) & nmol/L & 50.0 & PRIOR & Used only when \texttt{feedback\_mode='saturating'}, which is not the manuscript default. \\

Dietary conversion rate & Serum 25(OH)D gain per \textmu g/day oral D3 intake & nmol/L per (\textmu g/day) & 0.07 & PRIOR & \todo{Order-of-magnitude match to reported steady-state responses of roughly 1.5--2.0 nmol/L per \textmu g/day \citep{cashman2008_2011}; the specific source(s), cited only as ``Cashman 2008/2011,'' are not fully identified in this repository (two candidate years suggest two papers, not one).} \\

Body-weight reference & Reference body weight for volumetric dilution & kg & 70.0 & PRIOR & Functional form (hyperbolic dilution in total body weight, not BMI) verified against \citet{drincic2012} ($n=686$); the 70 kg reference value itself is a population-average convention. Cross-checked directionally against oral-dosing data: \citet{ekwaru2014} ($n=17{,}614$) found obese participants need 2--3 times the oral dose of normal-weight participants for the same 25(OH)D response, close to but somewhat larger than this model's weight-ratio implied factor; \citet{walsh2017} reviews the same dilution mechanism and \citet{pannu2016} report a meta-regression in which weight loss raises serum 25(OH)D, consistent in direction. \\

$D_{\mathrm{SAT}}$ (previtamin-D3 photodegradation) & Session dose at which net previtamin-D3 production saturates & SED & 7.5 & SINGLE-SOURCE PRIOR & \todo{\citet{vandijk2016}'s abstract reports 7.5 SED for one measured solar spectrum and 1.25 SED for another, a sixfold spread; this parameter takes the larger value and has not been reconciled against the smaller one.} \\

Hydroxylation inflection & Biphasic 25-hydroxylation efficiency inflection (non-default mode) & nmol/L & 56.0 & SINGLE-SOURCE PRIOR & \citet{commsmed2025}; used only when \texttt{feedback\_mode='biphasic'}, not the manuscript default. \\

Hydroxylation Hill exponent & Steepness of the biphasic efficiency transition & dimensionless & 2.0 & UNCONFIRMED & \todo{The cited source reports only the inflection location, not the transition's steepness; this exponent is an assumed shape parameter, not a fitted or measured one.} \\

Personal dose geometry factor & Ambient-to-personal UV exposure ratio, standing posture & dimensionless & 0.35 & PRIOR & \todo{Central value from \citet{cockell2001} ($0.35\pm0.12$), consistent with the \citet{parisi2000,parisi2003} ranges (0.14--0.49); treated as a constant scalar though \citet{schmalwieser2018} recommends expressing it as a function of solar elevation. Every minutes-of-sun recommendation this model produces inherits this uncertainty directly.} \\

Cancer-risk cumulative weight & SCC-type index scaling on cumulative dose & dimensionless & 1.0 & UNCONFIRMED & Structural choice, not fitted to incidence data; see Section~\ref{sec:background} and Section~\ref{sec:limitations}. \\

Cancer-risk intermittency exponent & BCC/melanoma-type index exponent on per-session dose relative to MED & dimensionless & 3.0 & UNCONFIRMED & Order-of-magnitude choice consistent with \citet{elwood1997}'s qualitative direction, not a fitted value. \\

Cancer-risk FST susceptibility decay & Exponential decay of cancer-risk index with Fitzpatrick type & dimensionless & 0.45 & UNCONFIRMED & Chosen to match a roughly tenfold incidence gap between very fair and deeply pigmented populations; not fitted to a specific incidence dataset. \\

MED by Fitzpatrick type & Minimal erythemal dose & SED & I: 2.1, II: 2.6, III: 3.2, IV: 5.6, V: 7.5, VI: 15.2 & VERIFIED & \citet{shih2018}. \\

\end{longtable}
}

Table~\ref{tab:sources} lists the per-source vitamin-D-effective dose weighting priors used for $E_{\mathrm{source}}$. These are calibration inputs, not universal constants: Section~\ref{sec:background} describes evidence that a single conversion factor per source cannot hold across studies, and the codebase requires callers to select or override a source-specific value rather than defaulting silently.

{\footnotesize
\begin{longtable}{@{}p{2.6cm}p{1.9cm}p{1.9cm}p{1.9cm}p{5.5cm}@{}}
\caption{Source/spectrum calibration priors ($E_{\mathrm{source}}$), from \texttt{SOURCE\_CALIBRATION\_TABLE} in \texttt{vitd\_model.py}, transcribed 2026-07-31.}
\label{tab:sources}\\
\toprule
Source & Mean & Range low & Range high & Basis \\
\midrule
\endfirsthead
\multicolumn{5}{c}{\tablename\ \thetable\ (continued)}\\
\toprule
Source & Mean & Range low & Range high & Basis \\
\midrule
\endhead
\bottomrule
\endfoot
Solar, summer & 1.89 & 1.79 & 2.00 & \citet{webb2021,grigalavicius2015} \\
Solar, broadband (reference) & 1.00 & 0.90 & 1.10 & Reference value; default for the model \\
TL-01 narrowband UVB & 2.44 & 2.24 & 2.63 & Midpoint of \citet{webb2021} (2.24) and \citet{grigalavicius2015} (2.63) \\
TL-12 broadband UVB & 1.24 & 1.23 & 1.25 & \citet{webb2021,grigalavicius2015} agree closely \\
UV6 broadband & 1.54 & 1.54 & 1.55 & \citet{webb2021,grigalavicius2015} agree closely \\
Arimed B & 1.83 & 1.83 & 1.83 & \citet{webb2021} only \\
Sunbed UVB & 1.18 & 1.18 & 1.18 & \citet{grigalavicius2015} Table 2; \citet{webb2021} reports a higher value (1.85) for a different tanning-lamp product and should not be substituted here \\
\end{longtable}
}

% =============================================================================
\section{Simulation setup and validation status}
\label{sec:validation}
% =============================================================================

\todo{This section is a stub. The repository's automated validation suite currently reports 9 of 10 required checks passing and 2 of 2 warning-only checks passing (regenerable with \texttt{run\_all.py}); the one failing required check is a NHANES seasonal-amplitude comparison, run on the RIA/LC-MS/MS-harmonized serum values documented in \citet{nhanes2015analytical}, confounded by unweighted survey means and season/behavior/latitude effects that the check does not control for; its failure mode (an inverted ratio) is treated as a diagnostic of feedback-term strength rather than archived as a dead end. A calibration target, \citet{bogh2012bjd}'s narrowband-UVB deficiency trial, is matched closely (model +56.0 nmol/L vs.\ observed +55.8 nmol/L) after a dose-unit correction rather than a parameter change. The model's only genuine held-out benchmark, not used to fit any parameter, is \citet{webbkift2018}: the model predicts a 1.40-fold skin type V to skin type II dose ratio against a reported 2.5- to 3.0-fold ratio, a clear failure that Section~\ref{sec:limitations} discusses. Beyond calibration and validation, three diagnostic scripts (\texttt{scripts/metric\_redundancy.py}, \texttt{scripts/parameter\_optima.py}, \texttt{scripts/melanin\_sc\_coupling.py}) probe the model's own internal consistency, whether its scheduling metrics agree with each other, how the optimal interval moves under fixed-budget versus per-SED framings, and whether melanin and stratum corneum can be told apart from each other; Section~\ref{sec:model} and Section~\ref{sec:limitations} report their results in full because they bear directly on what this model can and cannot claim, not because they are calibration or held-out checks in the sense above. A complete simulation-setup section (numerical integration scheme, schedule generators, and reproducibility details) remains to be written for the next revision of this manuscript.}

% =============================================================================
\section{Results}
\label{sec:results}
% =============================================================================

\todo{This section is a stub. Quantitative scheduling results (fixed-session-budget spacing, deficiency-versus-cancer-risk frontiers, and minutes-of-sun tables across Fitzpatrick types) exist in the accompanying codebase's generated output (\texttt{RESULTS.md}, regenerated by \texttt{run\_all.py}) but have not yet been selected, tabulated, and written up for this manuscript. Any number quoted from an out-of-date copy of that file should be re-verified by re-running the pipeline rather than assumed current.}

% =============================================================================
\section{Discussion}
\label{sec:discussion}
% =============================================================================

\todo{This section is a stub. A full discussion belongs here once Section~\ref{sec:results} is written, and should weigh the scheduling recommendations against the provisional status of $D_{\max}$ and $K_{\mathrm{ATTEN}}$ documented in Section~\ref{sec:limitations} rather than presenting them as settled.}

% =============================================================================
\section{Limitations and assumptions}
\label{sec:limitations}
% =============================================================================

This section states plainly what the model does not yet establish. Every item below is tracked in the source repository's own documentation (\texttt{MODEL\_AND\_EVIDENCE.md}) and is cross-referenced here so a reader does not have to trust a summary; the repository is the living version of this record and should be checked directly before any of these numbers are quoted elsewhere.

\paragraph{The homeostatic ceiling $D_{\max}$ has no confirmed source and drives the headline scheduling result.} \todo{No primary source in this repository supports $D_{\max}=150$ nmol/L. An earlier project revision attributed it to \citet{bogh2010jid}, whose reported ceiling was 116 nmol/L, not 150; a PMID once associated with this parameter belongs to an unrelated obesity study \citep{doumatey2010}.} This would matter less if $D_{\max}$ were a minor parameter. It is not. Ablating melanin attenuation entirely ($K_{\mathrm{ATTEN}}=0$) changes the model's ranking of exposure intervals by about 0.3\%; ablating melanin and stratum corneum attenuation together leaves the ranking intact. The feedback term $\max(0,1-D/D_{\max})$, not photoadaptation, produces the interval ordering that earlier phases of this project treated as a photoadaptation finding.

\paragraph{No objective checked so far supports a specific interval recommendation, and the most commonly quoted one actively points the wrong way.} The metric \texttt{gross\_synth\_per\_SED} (average vitamin D synthesized per unit UV dose delivered) rises monotonically as dosing frequency approaches zero, because sparse dosing keeps mean serum D low and each pulse meets less feedback inhibition. \todo{A rank-correlation sweep across a 1- to 60-day interval grid (\texttt{scripts/metric\_redundancy.py}) quantifies how badly this metric serves as a proxy for outcomes that matter: its ranking of schedules is close to the exact opposite of the ranking produced by mean 25(OH)D achieved over the horizon (Spearman $\rho \approx -0.999$) and of the ranking produced by avoiding time-integrated deficiency below 50 nmol/L (also $\rho \approx -0.999$ in the deficiency-avoidance direction). At the schedules this metric ranks best, 21- to 37-day intervals for skin type V in earlier project runs, final serum 25(OH)D sits below the starting baseline: a net loss of vitamin D, scored as maximally efficient. The same sweep found several metric pairs in the codebase rank-identical or rank-inverted (mean 25(OH)D and time-integrated deficiency below 50 nmol/L agree exactly, $\rho = +1.000$; area-under-the-curve per SED and final 25(OH)D disagree exactly, $\rho = -1.000$), meaning some of the eight scoring functions this project carries are not independent evidence for anything.}

An earlier draft of this manuscript, following this project's own prior documentation, reported that \texttt{net\_final\_D\_change\_per\_SED} removes the gross-efficiency artifact and locates a genuine interior optimum near three days. \todo{That claim needs correcting. \texttt{net\_final\_D\_change\_per\_SED} is still a per-SED ratio: cumulative UV dose is free to vary across the schedules it compares, the same property that lets \texttt{gross\_synth\_per\_SED} go wrong. \texttt{scripts/parameter\_optima.py} computed the optimal interval under a framing immune to this problem, cumulative UV dose held exactly fixed and split across sessions so that no schedule can win by dosing less, across the same fourteen limiting parameters described in Section~\ref{sec:model}, each swept over its full documented prior range and jointly by Monte Carlo. Under that framing, the optimal interval is one day, the finest spacing tested, for every parameter at every tested value, and the joint Monte Carlo gives a 90\% interval of exactly $[1,1]$ days. The three-day figure appears only under the per-SED ratio framing, and only two of the fourteen parameters tested, exposed body-surface-area fraction and the synthesis coefficient $A$, move that per-SED optimum at all, and both are dose-magnitude parameters rather than photoadaptation ones. This model currently supports no framing-independent interior optimal interval. The one-day result under the fixed-budget framing is itself a property of this model's concave, logarithmic dose-response curve, which rewards spreading a fixed total dose as finely as the sampling resolution allows; it is a computed property of this model's mathematics, not an independent finding about human photoadaptation, and should not be read as a clinical recommendation to dose daily either. No number of days between sessions should be quoted from this project as a settled result.}

\paragraph{The melanin attenuation coefficient rests on one seven-person study, and its most-quoted number is not what it is usually taken to mean.} $K_{\mathrm{ATTEN}}$ derives from \citet{matsuoka1990jlcm}, the only within-subject tanned-versus-untanned vitamin-D-synthesis study located for this project, abstract only, group means, $n=7$. \citet{neville2020} independently confirms the field still lacks replication of this result. \todo{The widely cited figure from this source, a 51\% reduction in vitamin D synthesis in tanned versus untanned skin, is confounded: the tanned arm also started with 64\% higher baseline 25(OH)D (92.1 vs.\ 56.2 nmol/L), and the study's own authors describe this as a ``resetting of the vitamin D3 synthetic mechanism,'' language that names feedback, not melanin. Deconvolving the baseline-feedback contribution from the melanin contribution leaves a melanin-only attenuation of 7 to 18\%, not 51\%, depending on which feedback functional form is assumed; the implied $K_{\mathrm{ATTEN}}$ range under that deconvolution is 0.34 to 1.00, and the model's current value of 0.37 falls inside that range. Do not cite the 51\% figure as a melanin effect anywhere it is repeated from this paper.}

\paragraph{Three parameter pairs cannot be separately identified from the data available to this project.} Sensitivity analysis shows that more than 99\% of the response-surface variance for each of three pairs, $K_{\mathrm{ATTEN}}$ and $K_{SC}$, $\alpha_{SC}$ and $SC_{\max}$, and $K_{\mathrm{const}}$ and $K_{\mathrm{ATTEN}}$, lies along a single ridge: only the combination of each pair is constrained by current data, not the individual members. Reporting either member of a pair as independently calibrated overstates what the evidence supports. Resolving this requires individual-level trajectory data (melanin index and stratum corneum thickness measured on the same participants over time), which is not currently available; \citet{deboer2025}'s data-availability statement indicates it may be obtainable on request.

\paragraph{Melanin attenuation is nearly inert at realistic doses, and most of the modeled skin-type gap comes from constitutive pigmentation and stratum corneum, not tanning.} \todo{Sweeping $K_{\mathrm{ATTEN}}$ across its full plausible range (0 to 1) at a realistic dose, 1 SED every three days for skin type V over a year, changes mean 25(OH)D by 0.16 nmol/L; the melanin index itself moves from 0.780 to 0.791 over the same year (\texttt{scripts/melanin\_sc\_coupling.py}). A supplementary check found the effect reaches a still-modest 3.2 nmol/L only at 4 SED per day, a cumulative 1460 SED per year that is not a realistic exposure at any latitude. An ablation of the steady-state 25(OH)D gap this model predicts between skin types II and V (7.44 nmol/L as currently configured) attributes roughly 76\% of that gap to constitutive pigmentation, roughly 18\% to stratum corneum attenuation, and approximately 0\% to facultative tanning; removing the tanning channel entirely widens the modeled gap by 1.7\% rather than narrowing it (\texttt{scripts/melanin\_sc\_coupling.py}). Consistent with the identifiability ridge described above, $K_{\mathrm{ATTEN}}$ and $K_{SC}$ place 99.94\% of their combined effect on steady-state 25(OH)D along one direction, but asymmetrically: $K_{SC}$'s marginal effect on steady-state D across its prior range (6.57 nmol/L) is roughly forty times $K_{\mathrm{ATTEN}}$'s (0.16 nmol/L). Read together, these are computed properties of this model as currently configured, not independent findings about human photoadaptation. They say plainly, in the model's own terms, that facultative tanning barely shapes anything this model computes at realistic doses, and that this manuscript's framing of tanning as a central mechanism (Section~\ref{sec:background}) is carrying far more rhetorical weight than quantitative weight in the model's actual behavior.}

\paragraph{The stratum-corneum optical-attenuation mechanism may not be the mechanism actually occurring.} \todo{$K_{SC}$ models stratum corneum thickening as added optical path length attenuating UV before it reaches the synthesis-relevant epidermal layer. In \citet{deboer2025}'s data, epidermal thickness barely changed (108.6 to 109.4 \textmu m, $P=0.718$) while stratum corneum thickness swung by roughly 26\% in each direction over nine weeks; the authors attribute this to edema, not acanthosis. If the underlying change is fluid retention rather than additional corneocyte layers, the optical-attenuation interpretation of $K_{SC}$ is mechanistically wrong even where the fitted number happens to improve agreement with data. Combined with the identifiability problem above, this argues for treating stratum corneum and melanin as one empirical photoadaptation channel rather than defending two separately mechanistic ones, though this manuscript has not made that structural change.}

\paragraph{The model fails its only genuine held-out benchmark.} \citet{webbkift2018} is the sole study in this project not used to fit or calibrate any parameter. Against its reported 2.5- to 3.0-fold dose penalty for skin type V relative to skin type II, the model predicts 1.40-fold, a clear miss. A looser check on required daily minutes of sun passes within a $\pm40\%$ band, but the dose-ratio comparison is the cleaner test (ambient UV index and personal-dose geometry cancel out of a ratio in a way they do not cancel out of an absolute minutes prediction), and it fails. The likely explanation is that $K_{\mathrm{const}}$, this model's constitutive-pigmentation penalty, is calibrated to \citet{young2020}'s per-unit-dose ratio under matched-BSA lamp exposure (roughly 1.3- to 1.4-fold), which is a different quantity from \citet{webbkift2018}'s whole-body naturalistic dose penalty (2.5- to 3.0-fold). Every real-world sun-exposure recommendation this model produces for darker skin types should be read with this under-prediction in mind, and probably underestimates the true skin-type penalty.

\paragraph{The Fitzpatrick type V and VI melanin mapping is extrapolated, and pigmentation is a weak individual predictor regardless.} \todo{$M_{\mathrm{baseline}}$ for skin types V (0.78) and VI (1.00) has no direct published melanin-index validation and is extrapolated from an $n=44$ trial dataset anchored at lighter skin types. Two papers give partial, instrument-specific anchors for dark skin without resolving this: \citet{wilkes2015} report that a Mexameter melanin index above 750 maps to Fitzpatrick type VI in an African population, and \citet{alaluf2001} quantify melanin content directly in type V and VI skin by HPLC, but neither gives a matched untanned type V versus type VI baseline on the reflectance scale this model uses.} More fundamentally, \citet{datta2019} found that measured skin pigmentation explains only 1.9 to 2.6\% of individual variance in UVB-induced vitamin D response; pigment-associated genetic variants, baseline serum D, and body-surface-area exposure dominate. Every skin-type-indexed parameter in this model should be read as a group-level phenotype marker, not a precise predictor for any one person.

\paragraph{Skin cancer risk is a qualitative ranking device, not a validated quantitative model, by design.} As stated in Section~\ref{sec:background}, this project deliberately deprioritizes quantitative skin-cancer risk modeling. The implemented indices use structurally defensible but unfitted exponents, and the papers that would allow fitting them, \citet{eadie2021} \todo{not in this repository}, \citet{armstrongkricker2001} \todo{not in this repository}, and a 2023 PLOS ONE study \citep{plosone2023} \todo{not in this repository, no author recorded}, are absent from the source corpus. Treat any cancer-risk number this model produces as a same-total-dose ranking device only.

\paragraph{Additional open items.} The personal-dose geometry factor that converts ambient UV index into UV dose actually received by a standing person (0.35) is a constant scalar; \citet{schmalwieser2018} recommends expressing this ratio as a function of solar elevation, and until that is implemented every minutes-of-sun recommendation carries at least $\pm40\%$ geometric uncertainty on top of its other sources of error. \todo{The TL-01 lamp-to-erythemal-SED conversion factor (5.5, used to correct the \citet{bogh2012bjd} calibration check) comes from a single cabinet measurement \citep{rodriguezgranados2017} and is cabinet-specific: \citet{lloyd2004} found TL-01 calibration factors varying by $\pm60\%$, up to 2.7-fold, across seven expert UK medical physics departments, so any TL-01-derived absolute prediction in this model carries roughly twofold dose uncertainty on top of everything else.} The previtamin-D3 photodegradation saturation dose ($D_{\mathrm{SAT}}=7.5$ SED) rests on a single abstract reporting a sixfold spread (7.5 SED for one measured spectrum, 1.25 SED for another) that has not been reconciled. This model does not include a genetic pathway component; \citet{datta2019,datta2016} both show pigment-associated genetic variants explain substantially more individual variance than measured pigmentation, and no dataset linking genotype, UV exposure, and 25(OH)D response was available to this project. No naturalistic solar-exposure validation dataset exists for skin types V and VI in this repository; every skin type V and VI data point used here comes from a lamp source (TL-01 or Arimed B), not natural sunlight.

% =============================================================================
\section{Conclusion}
\label{sec:conclusion}
% =============================================================================

\todo{This section is a stub pending Section~\ref{sec:results} and Section~\ref{sec:discussion}.} The contribution drafted in this manuscript is structural: a three-state model that lets pigmentation and barrier state evolve with an exposure schedule, three decision-framed objectives in place of a single dose-response slope, and an explicit, source-traced evidence-status tag on every parameter. This manuscript deliberately stops short of recommending a specific inter-session interval: Section~\ref{sec:limitations} shows that the model's own optimal spacing depends entirely on which objective is asked, collapses to the finest spacing tested under the one framing immune to ratio artifacts, and that melanin attenuation, the mechanism motivating this project, moves the model's output by well under one nmol/L at doses anyone would actually receive. The scientific content that would justify a specific scheduling recommendation, chiefly a defensible replacement for $D_{\max}$, a resolution of the constitutive-pigmentation calibration gap, and an objective that does not degenerate under either dosing extreme, remains future work, and the user of this model should treat every quantitative output as provisional until those items are resolved.

% =============================================================================
% End matter
% =============================================================================

\section*{CRediT authorship contribution statement}
\todo{To be completed. Single-author project at time of writing; CRediT roles (conceptualization, methodology, software, writing) to be assigned explicitly before submission.}

\section*{Declaration of competing interest}
\todo{To be completed by the author before submission.}

\section*{Data availability}
The model code, validation suite, and evidence documentation underlying this manuscript are available in the project repository. \todo{Repository URL / DOI / archival deposit (e.g.\ Zenodo) to be added before submission.} NHANES, trial-CSV, and PDF sources referenced throughout are third-party data; see the repository's \texttt{data\_provenance.yaml} for per-dataset provenance and access terms.

\section*{Acknowledgments}
\todo{To be completed.}

\bibliographystyle{elsarticle-harv}
\bibliography{references}

\end{document}
